\section{Segunda iteración de la Fase 5: Formato del Cartucho}
\label{sec:F5I2}

El análisis de la iteración anterior dejó el inventario cerrado: en el GWC viaja todo lo que los dos motores leen en un \textit{tick}, y nada más. Esta iteración define la estructura exacta del fichero que se va a generar intentando optimizarlo al máximo.

Su definición exacta se incluye en el apéndice  \ref{sec:Formatos} al final de este documento.

\subsection{Generación determinista del cartucho}
\label{sec:F5I2Def}

La generación es siempre \textbf{determinista}: el mismo GWI debe
producir siempre el mismo GWC byte a byte, con el mismo CRC. El generador es equivalente a un compilador, con
el mismo origen siempre produce el mismo resultado.

Para asegurar esto, se añade en el Cartucho GWI un campo denominado \textbf{codigoGWC} con el contenido del campo Checksum del cartucho GWC junto a sus marcas de inicio y fin (formato \#FFFF\#), que coincide cuando ambos cartuchos se han generado a partir de los mismos datos exactamente.

\subsection{Posibles mejoras}
\label{sec:F5I2Mejoras}

Para reducir más el tamaño del cartucho, se pueden usar en algunos lugares ciertas mejoras, pero en general cambiarían espacio por tiempo de procesamiento, lo que no se ve necesario en este momento y no se abordarán en esta fase.

La principal es que si el cartucho tiene pocos segmentos se pueden reducir el tamaño de los punteros: se podría usar el campo con el contador de segmentos para decidir el tamaño real de los punteros, que seguramente se puede reducir a solo 2 bytes en muchos juegos.

Los registros de segmento usan campos que no se emplean en todos los tipos; cada registro reserva siempre espacio para 7 barras, 4 dígitos, 4 direcciones, etc. Usando registros de longitud variable dependiendo del tipo se podrían llegar a reducir a segmentos de entre 14 bytes y 30 bytes, en lugar de usar siempre los 32 actuales, lo que puede llegar a representar hasta 1~KB menos en juegos con 80 segmentos.

